\section{Generation and validation}
\label{sec:arch-generation}

The shared computational core needs several interfaces: typed C++ calls,
PostgreSQL functions, configuration settings, and benchmark adapters. Keeping
these interfaces consistent is part of the architecture. Generation removes
some repeated declarations, while assigning compilation work to modules
limits its repetition.

\subsection{Deriving the PostgreSQL boundary}
\label{sec:arch-sql-generation}

A PostgreSQL routine repeats one contract in several places. Installation SQL
declares its argument types, result and volatility. A C entry point must expose
the corresponding symbol, decode those arguments and encode that result.
BioSeqDB derives these artifacts from annotated C++ declarations. The following
existing annotation fixture illustrates the API. Its declaration supplies the
name and C++ types, while attributes supply SQL metadata:

\noindent\begin{minipage}{\linewidth}
\begin{lstlisting}[style=bioseq-cpp26,basicstyle=\footnotesize\ttfamily,
  caption={A typed declaration with SQL metadata.},
  label={lst:arch-sql-annotation}]
[[bioseq_sql::function({
    .volatility = bioseq::sql::volatility::immutable,
})]]
inline constexpr std::int32_t identity(
    [[bioseq_sql::param({.sql_type = bioseq::sql::pg::integer})]]
    std::int32_t value) {
    return value;
}
\end{lstlisting}
\end{minipage}

The generator uses Clang's C++ LibTooling and abstract syntax tree (AST) interfaces through
\texttt{libclang-cpp}. Compiler actions inspect registered module owners under
an explicit compilation profile~\autocite{clang_libtooling,clang_tooling_interfaces}.
AST facts describe declarations and types. Compiler token ranges associate
annotations and documentation with those declarations.
BioSeqDB then parses the annotation options as a restricted language and
validates their admitted forms. Unknown options, conflicting annotations,
duplicate defaults and unsupported C++/SQL conversions stop generation.
Volatility and parallel safety remain author-supplied contracts: the generator
does not prove them from an arbitrary function body.

\begin{figure}[htbp]
\centering
\begin{tikzpicture}[
  font=\small,>=Latex,
  stage/.style={draw,rounded corners,align=center,fill=blue!6,
    minimum height=1.02cm,text width=2.9cm},
  sqlobject/.style={draw,rounded corners,fill=white,align=center,
    font=\footnotesize,minimum height=.43cm,inner sep=3pt},
  artifact/.style={draw,rounded corners,align=center,fill=green!8,
    minimum height=.72cm,text width=2.7cm}]
  \node[stage] (source) at (0,0) {Annotated declaration\\C++ types and SQL metadata};
  \node[stage] (clang) at (3.65,0) {Clang AST and tokens\\active modules and profile};
  \node[draw,rounded corners,fill=blue!6,
    minimum width=5.4cm,minimum height=1.45cm] (graph) at (8.35,0) {};
  \node at (8.35,.46) {SQL object dependencies};
  \node[sqlobject,text width=1.12cm] (schema) at (6.55,-.20) {schema};
  \node[sqlobject,text width=1.0cm] (type) at (8.35,-.20) {type};
  \node[sqlobject,text width=1.32cm] (function) at (10.15,-.20) {function};
  \draw[->] (schema) -- (type);
  \draw[->] (type) -- (function);
  \draw[->,thick] (source) -- (clang);
  \draw[->,thick] (clang) -- (graph);
  \node[artifact] (sql) at (-.1,-2.0) {Installation SQL};
  \node[artifact] (fmgr) at (3.2,-2.0) {C ABI wrappers};
  \node[artifact] (support) at (6.5,-2.0) {Typed support modules};
  \node[artifact] (manifest) at (9.8,-2.0) {Build metadata};
  \draw[thick] (graph.south) -- (8.35,-1.12);
  \draw[thick] (-.1,-1.12) -- (9.8,-1.12);
  \foreach \target in {sql,fmgr,support,manifest}
    \draw[->,thick] (\target.north |- 0,-1.12) -- (\target.north);
\end{tikzpicture}
\caption{One checked interface produces related artifacts. The small dependency
chain is illustrative. Typed support includes row codecs, configuration
bindings and catalog handles where the declarations require them.}
\label{fig:arch-sql-generation}
\end{figure}

The graph records dependencies among declared SQL objects, each with a stable
identity.
Tarjan's algorithm identifies strongly connected components, after which a
stable topological traversal orders the admitted component
graph~\autocite{tarjan_1972_depth_first_search}. Unsupported cycles are rejected.
Base types additionally need statement-level ordering between their shell
declaration, input/output functions and completed definition.

Generated support code must exist before its consumers compile. The build
therefore produces row codecs, configuration bindings and catalog handles
before final SQL and wrappers. Final generation checks these intermediate
outputs against the active declarations. Separate build profiles keep
test-only declarations out of the production interface.

The generated C ABI also centralizes exception translation. C++ exceptions
must not escape into PostgreSQL, while a PostgreSQL error can jump past C++
destructors~\autocite{postgresql17_cpp_extensibility}. The boundary preserves captured
database diagnostics and leaves C++ exception scopes before raising the SQL
error. Calls into PostgreSQL below that boundary still require an explicit
bridge or a region without live nontrivial C++ objects.

Language reflection is a related alternative, but it is not the mechanism
used here. Reflection was adopted for C++26~\autocite{wg21_n5015_reflection}.
Upstream Clang still lists it as unsupported~\autocite{clang_cxx_status}. BioSeqDB
therefore uses an external compiler pass, including source locations,
preprocessing state and module provenance, together with its own annotation
validation.

\subsection{Controlling specialization cost}
\label{sec:arch-specialization-cost}

Templates specialize kernels for concrete alphabets, storage layouts and
algorithm choices. Independent parameters multiply the potential number of
instantiations (Figure~\ref{fig:arch-template-product}). A parameter can create
new compiler work even when it does not change a downstream algorithm.
Link-time removal of equivalent code cannot recover all the earlier parsing,
instantiation and optimization work.

\begin{figure}[htbp]
\centering
\begin{tikzpicture}[
  font=\small,>=Latex,
  cell/.style={draw,fill=blue!7,minimum width=.63cm,
    minimum height=.50cm,inner sep=0pt,font=\footnotesize}]
  \foreach \implementation/\offset in {1/0,2/4.2,3/8.4} {
    \begin{scope}[xshift=\offset cm]
      \node at (1.02,.50) {Implementation $I_{\implementation}$};
      \foreach \column/\width in {0/8,1/16,2/32,3/64}
        \node[font=\footnotesize] at (.68*\column,-.12) {\width};
      \foreach \row/\alphabet in {0/1,1/2} {
        \node[font=\footnotesize] at (-.53,-.70-.57*\row) {$A_{\alphabet}$};
        \foreach \column in {0,1,2,3}
          \node[cell] at (.68*\column,-.70-.57*\row) {};
      }
      \node[cell,fill=orange!25] at (1.36,-.70) {$f_{\implementation}$};
      \draw[thick] (-.315,-1.54) -- (-.315,-1.66)
        -- (2.355,-1.66) -- (2.355,-1.54);
    \end{scope}
  }
  \node[draw,rounded corners,fill=green!8,align=center,
    text width=10.7cm,minimum height=.72cm] (helper) at (5.22,-2.48)
    {Shared helper body independent of these dimensions};
  \foreach \x in {1.02,5.22,9.42}
    \draw[->,thick] (\x,-1.66) -- (helper.north -| \x,0);
\end{tikzpicture}
\caption{Illustrative specialization product: two alphabets, four storage
widths in bits, and three implementations produce 24 cells. Each cell denotes
one possible function body. Highlighted cells keep alphabet and width fixed
while changing implementation. This is not a production support matrix.}
\label{fig:arch-template-product}
\end{figure}

BioSeqDB limits this cost by selecting length buckets after fixing the k-mer
size, passing already selected input representations into scheduled work, and assigning
expensive specializations to concrete module owners. Shared computational
modules can be compiled once for compatible consumers. An ordinary benchmark
and the extension can consequently reuse the same core without each owning
its compilation again. This does not eliminate specialization: it gives the
build graph explicit work units and avoids repeating irrelevant choices.

The cost has been substantial during development. Recorded May 2026 runs
included a 32-minute debug benchmark build and PostgreSQL translation-unit
compilations that exceeded a 50-minute timeout. These are historical
observations, not current build timings or measurements of SQL generation
alone. Later investigations also found expensive chains of imported modules
behind small function bodies. Controlling both template dimensions and module
ownership is therefore necessary.

\subsection{Benchmark specifications and validation}
\label{sec:arch-benchmark-specifications}

Benchmarking uses a generic framework and adapters for individual workloads.
The framework owns configuration validation, matrix expansion, scheduling,
measurement and report generation. An adapter defines the experiment: how to
prepare its input, invoke an algorithm or SQL statement, evaluate the result
and release temporary resources. The same orchestration serves in-process
alignment, sequence input and output, and PostgreSQL experiments.

A JSON specification declares named inputs and algorithm parameter schemas.
Within each benchmark, \texttt{matrix.for} assigns repetition counts to inputs,
while \texttt{matrix.over} lists algorithms and parameter values. The framework
expands the combinations separately for each algorithm and fills omitted
parameters from declared defaults. Schema and reference checks reject
inconsistent specifications before execution. JSON selects implementations
already registered by the adapters.

\noindent\begin{minipage}{\linewidth}
\begin{lstlisting}[style=bioseq-arch-code,language={},
  caption={Input and algorithm selection fields from an alignment benchmark specification.},
  label={lst:arch-benchmark-matrix}]
"matrix": {
  "for": {"RMQ_REPRO_LONG_CCS": 100},
  "over": {
    "bioseq_ksw2_rmq_scalar_mm": {},
    "minimap2_rmq": {}
  }
}
\end{lstlisting}
\end{minipage}

This excerpt selects 100 repetitions of each registered alignment route on
the same named input. The empty objects use the declared parameter defaults.
Input definitions, metric selection and other run settings are omitted.

The measurement boundary is explicit. The framework times the execution
callback, then passes its input and output to the evaluator. This permits
reference computations outside that timer. However, adapters may also check
results while constructing their output, so elapsed time must be interpreted
together with the measured region. Alignment experiments through PostgreSQL
check expected matches, coordinates within the configured tolerance and
strand. Selected cases additionally validate CIGAR replay. These checks
complement authored integration tests of SQL behavior.

Reports retain metric summaries and rule outcomes. Adapters can additionally
emit durable raw samples. Rules compare selected configurations within one
run, while regression checks use eligible historical captures. Eligibility
requirements determine which runs can enter a historical comparison.
The configured matrix still determines coverage: a successful timing run
does not establish correctness for every query, setting or executor path.
